\chapter{Tailor a benchmark towards a particular machine}

Codes and machines are not independent of each other: 
Even the best and most flexible codes on earth will not be able to exploit a given supercomputer if they are not translated with the right compilers, are using the correct system settings, and so forth.
In the following sections, we evaluate which settings have to be chosen to make a code perform on a given system.
The organisation of the sections follows the rubrics in Chapter~\ref{chapter:high-level}.


In practice, it is one of the first steps of a system specialist to tweak and tailor the setting around a code to the given target machine.
In an ideal case, this process is guided by systematic performance assessment from Section~\ref{chapter:high-level}.   
In turn, that implies that we have to implement an iterative process, where each tailoring is followed by a rerun of the high-level assessment.


